% =============================================================================
% APPENDIX FFF: METHOD-REGISTRY: Model Archive and Organizational Learning
% =============================================================================
% Module: BCM2_METHOD_REGISTRY
% Category: METHOD
% Language: English
% Version: 1.0 (January 2026)
% Status: Draft
%
% This appendix defines the Model Registry---a structured archive of all
% behavioral models designed using the EBF framework. The Registry enables
% organizational learning, provides options for the 3+1 Choice Architecture,
% and maintains audit trails for model-based recommendations.
%
% =============================================================================

% =============================================================================
% CROSS-REFERENCE STRUCTURE
% =============================================================================
%
% CHAPTER LINKAGE (Required):
% - Primary Chapter: Chapter 4: From Framework to Model
% - Secondary Chapters: Chapter 18 (Meta-Theory), Chapter 10 (Applications)
%
% DEPENDENCIES (what this appendix REQUIRES):
% - Appendix UNMAPPED_DSN (METHOD-DESIGN): Defines model structure and workflow
% - Appendix UNMAPPED_EST (CORE-WHERE): Parameter source documentation
% - All DOMAIN- appendices: Domain-specific model attributes
%
% DEPENDENTS (what REQUIRES this appendix):
% - Appendix UNMAPPED_DSN (METHOD-DESIGN): Draws 3+1 options from Registry
% - All PREDICT- appendices: Models stored here
% - Quality Assurance processes: Validation tracking
%
% =============================================================================

\begin{tcolorbox}[colback=purple!5!white, colframe=purple!75!black,
    title=Chapter Linkage]

\textbf{Primary Chapter:} Chapter 4: From Framework to Model
\begin{itemize}[nosep]
\item Section 4.4: Model Storage $\leftrightarrow$ Section \ref{sec:fff-schema}
\item Section 4.5: Organizational Learning $\leftrightarrow$ Section \ref{sec:fff-learning}
\end{itemize}

\textbf{Secondary Chapters:}
\begin{itemize}[nosep]
\item Chapter 18: Meta-Theory --- Registry as knowledge management
\item Chapter 10: Applications --- domain-specific model retrieval
\end{itemize}

\textbf{Reading Guide:}
\begin{itemize}[nosep]
\item For model design process: Read Appendix UNMAPPED_DSN first
\item For storage structure: This appendix provides complete specification
\item For parameter sources: See Appendix UNMAPPED_EST (CORE-WHERE)
\end{itemize}

\end{tcolorbox}

\chapter{METHOD-REGISTRY: Model Archive and Organizational Learning}
\label{app:method-registry}

\begin{abstract}
This appendix defines the Model Registry---a structured archive storing all
behavioral models designed from the Evidence-Based Framework (EBF). The Registry
serves three functions: (1) enabling the 3+1 Choice Architecture through
similarity search, (2) maintaining audit trails for model-based recommendations,
and (3) facilitating organizational learning across projects. The Registry schema
captures model metadata, inputs, specifications, outputs, and validation status.
\end{abstract}

\begin{tcolorbox}[colback=blue!5!white, colframe=blue!75!black,
    title=Quick Reference: Model Registry]
\textbf{Appendix:} REG (METHOD)


\textbf{Core Question:} How do we store, retrieve, and learn from behavioral models?

\textbf{Key Structure:}
\begin{equation}
\mathcal{R} = \{M_1, M_2, \ldots, M_n\} \quad \text{where} \quad M_i = (\text{ID}, \text{Meta}, \text{Input}, \text{Spec}, \text{Output}, \text{Val})
\end{equation}

\textbf{Key Concepts:}
\begin{itemize}[nosep]
\item \textbf{Model Schema}: Standardized structure for model documentation
\item \textbf{Similarity Search}: Finding comparable models for 3+1 options
\item \textbf{Validation Status}: Tracking prediction accuracy over time
\item \textbf{Audit Trail}: Traceability of model-based recommendations
\end{itemize}

\textbf{Cross-References:} EEE (Design Workflow), BBB (Parameters), DOMAIN- (Applications)
\end{tcolorbox}

% -----------------------------------------------------------------------------
% APPENDIX SCOPE BOX
% -----------------------------------------------------------------------------

\begin{tcolorbox}[
    colback=orange!5!white,
    colframe=orange!75!black,
    title={\textbf{Appendix Scope: Ziel / In-Scope / Out-of-Scope / Constraints}},
    fonttitle=\bfseries
]

\textbf{Ziel (Objective):}
\begin{itemize}[nosep]
    \item How do we systematically store, retrieve, and learn from behavioral
models to enable organizational memory and continuous improvement?
\end{itemize}

\textbf{In-Scope:}
\begin{itemize}[nosep]
    \item The Fundamental Question
    \item Model Schema
    \item Similarity Search
    \item Organizational Learning
    \item Audit Trail
\end{itemize}

\textbf{Out-of-Scope (delegiert an andere Appendices):}
\begin{itemize}[nosep]
    \item METHOD-DESIGN $\rightarrow$ Appendix UNMAPPED_DSN
    \item CORE-WHERE $\rightarrow$ Appendix UNMAPPED_EST
    \item METHOD-DESIGN $\rightarrow$ Appendix UNMAPPED_DSN
    \item CORE-WHERE $\rightarrow$ Appendix UNMAPPED_EST
    \item Central Glossary $\rightarrow$ Appendix UNMAPPED_GLS
\end{itemize}

\textbf{Constraints (Anwendungsgrenzen):}
\begin{itemize}[nosep]
    \item [TODO: Define when this appendix does NOT apply]
    \item [TODO: Define required prerequisites]
    \item [TODO: Define known limitations]
\end{itemize}

\textbf{Lieferobjekte (Deliverables):}
\begin{itemize}[nosep]
    \item 5 Table(s)
    \item 5 Equation(s)
    \item 6 Axiom(s)
    \item Worked Example(s)
\end{itemize}

\end{tcolorbox}




% =============================================================================
%                    PART 2: CORE CONTENT
% =============================================================================

%%%%%%%%%%%%%%%%%%%%%%%%%%%%%%%%%%%%%%%%%%%%%%%%%%%%%%%%%%%%%%%%%%%%%%%%%%%%%%%
\section{The Fundamental Question}
\label{sec:fff-fundamental}
%%%%%%%%%%%%%%%%%%%%%%%%%%%%%%%%%%%%%%%%%%%%%%%%%%%%%%%%%%%%%%%%%%%%%%%%%%%%%%%

\subsection{Why This Matters}

Organizations design behavioral models for projects, apply them, and then...
forget them. Without systematic storage:

\begin{itemize}
\item Knowledge walks out the door when consultants leave
\item The same model is rebuilt from scratch for similar projects
\item No learning occurs from model successes or failures
\item Recommendations cannot be audited or justified retrospectively
\end{itemize}

\subsection{The Core Problem}

\begin{tcolorbox}[colback=red!5!white, colframe=red!75!black,
    title=The Fundamental Question]
\textbf{How do we systematically store, retrieve, and learn from behavioral
models to enable organizational memory and continuous improvement?}
\end{tcolorbox}

The Registry addresses this by providing:
\begin{enumerate}
\item \textbf{Structured Storage}: Every model has complete, standardized documentation
\item \textbf{Efficient Retrieval}: Similarity search finds relevant historical models
\item \textbf{Validation Tracking}: Model predictions are compared to outcomes
\item \textbf{Audit Trail}: Every recommendation can be traced to its model
\end{enumerate}


%%%%%%%%%%%%%%%%%%%%%%%%%%%%%%%%%%%%%%%%%%%%%%%%%%%%%%%%%%%%%%%%%%%%%%%%%%%%%%%
\section{Model Schema}
\label{sec:fff-schema}
%%%%%%%%%%%%%%%%%%%%%%%%%%%%%%%%%%%%%%%%%%%%%%%%%%%%%%%%%%%%%%%%%%%%%%%%%%%%%%%

Every model in the Registry follows a standardized schema with six components.

\subsection{Schema Overview}

\begin{figure}[htbp]
\centering
\begin{tikzpicture}[node distance=0.8cm,
    block/.style={draw, rectangle, minimum width=10cm, minimum height=0.8cm, align=left, font=\small}]

    \node[block, fill=blue!10] (id) {\textbf{1. IDENTIFICATION} \hfill Model\_ID, Designer, Created, Client};
    \node[block, fill=green!10, below=of id] (meta) {\textbf{2. METADATA} \hfill Gestaltungsanlass, Entry Point, Domain};
    \node[block, fill=yellow!10, below=of meta] (input) {\textbf{3. INPUT} \hfill Scope, Context ($\Psi$), Constraints};
    \node[block, fill=orange!10, below=of input] (spec) {\textbf{4. SPECIFICATION} \hfill Variables ($C$), Complementarity ($\gamma$), Parameters ($\Theta$), Form};
    \node[block, fill=red!10, below=of spec] (output) {\textbf{5. OUTPUT} \hfill Predictions, Statements, Recommendations};
    \node[block, fill=purple!10, below=of output] (val) {\textbf{6. VALIDATION} \hfill Status, Accuracy, Date, Evidence};

\end{tikzpicture}
\caption{Model Registry Schema: Six Components}
\label{fig:fff-schema}
\end{figure}

\subsection{Component 1: Identification}

\begin{tcolorbox}[colback=blue!5!white, colframe=blue!75!black,
    title=Schema: Identification]
\begin{verbatim}
identification:
  model_id: M-YYYY-NNN          # Unique identifier
  designer: string              # Name of model designer
  created: YYYY-MM-DD           # Creation date
  modified: YYYY-MM-DD          # Last modification
  client: string                # Client identifier (can be anonymized)
  project: string               # Project name/code
  version: X.Y                  # Model version
\end{verbatim}

\textbf{Naming Convention:}
\begin{itemize}[nosep]
\item M = Model prefix
\item YYYY = Year (e.g., 2026)
\item NNN = Sequential number within year (001, 002, ...)
\end{itemize}
\end{tcolorbox}

\subsection{Component 2: Metadata}

\begin{tcolorbox}[colback=green!5!white, colframe=green!75!black,
    title=Schema: Metadata]
\begin{verbatim}
metadata:
  gestaltungsanlass: string     # The design occasion (free text)
  gestaltungsanlass_type:       # Categorization
    - problem                   # Something wrong to fix
    - opportunity               # Something possible to pursue
    - question                  # Something unclear to understand
    - situation                 # Something to shape
  entry_point:                  # How model design started
    - theory_driven             # Started with EBF concept
    - practice_driven           # Started with situation
    - hybrid                    # Both constraints given
  domain: string                # Primary domain (e.g., "retail", "pension")
  tags: [string]                # Searchable keywords
\end{verbatim}
\end{tcolorbox}

\subsection{Component 3: Input}

\begin{tcolorbox}[colback=yellow!5!white, colframe=yellow!75!black,
    title=Schema: Input]
\begin{verbatim}
input:
  scope:
    description: string         # What behavior is modeled
    scope_type:                 # Category of behavior
      - decision                # Binary choice
      - continuous              # Amount, frequency
      - process                 # Sequence of behaviors
  context_psi:
    geographic: string          # Location (e.g., "Switzerland")
    demographic: string         # Target population
    temporal: string            # Time context
    channel: string             # Interaction channel
    regulatory: string          # Legal/regulatory environment
    cultural: string            # Cultural factors
    economic: string            # Economic conditions
    competitive: string         # Market structure
  constraints:
    budget: float               # Available budget (if applicable)
    timeline: string            # Project timeline
    data_available: [string]    # What data exists
\end{verbatim}
\end{tcolorbox}

\subsection{Component 4: Specification}

\begin{tcolorbox}[colback=orange!5!white, colframe=orange!75!black,
    title=Schema: Specification]
\begin{verbatim}
specification:
  variables_C:                  # Utility dimensions included
    - F                         # Financial
    - E                         # Emotional
    - P                         # Physical
    - S                         # Social
    - D                         # Developmental
    - X                         # Existential (if applicable)
  complementarity_gamma:
    interactions: [string]      # Which interactions (e.g., "F*E", "temporal")
    values: {string: float}     # Parameter values
  parameters_theta:
    values: {string: float}     # All parameter values
    sources: {string: string}   # Source for each parameter
    confidence: {string: float} # Confidence level [0,1]
  functional_form:
    type: string                # "linear", "cobb_douglas", "ces", "logistic", ...
    equation: string            # LaTeX equation
    justification: string       # Why this form was chosen
\end{verbatim}
\end{tcolorbox}

\subsection{Component 5: Output}

\begin{tcolorbox}[colback=red!5!white, colframe=red!75!black,
    title=Schema: Output]
\begin{verbatim}
output:
  predictions:                  # Testable predictions
    - id: P1
      statement: string         # The prediction
      type: point|comparative|boundary
      quantification: string    # Expected effect size
      testable: boolean         # Can be empirically tested
  statements:                   # What was communicated to client
    - date: YYYY-MM-DD
      audience: string          # Who received statement
      content: string           # What was said
      based_on: [P1, P2, ...]   # Which predictions supported it
  recommendations:              # Action recommendations
    - id: R1
      action: string            # Recommended action
      expected_impact: string   # Predicted impact
      based_on: [P1, P2, ...]   # Supporting predictions
\end{verbatim}
\end{tcolorbox}

\subsection{Component 6: Validation}

\begin{tcolorbox}[colback=purple!5!white, colframe=purple!75!black,
    title=Schema: Validation]
\begin{verbatim}
validation:
  status:                       # Current validation state
    - pending                   # Not yet validated
    - validated                 # Predictions tested and confirmed
    - partially_validated       # Some predictions confirmed
    - falsified                 # Predictions tested and rejected
    - superseded                # Replaced by newer model
  accuracy:                     # If validated
    overall: float              # Overall accuracy [0,1]
    by_prediction: {string: float}  # Per-prediction accuracy
  validation_date: YYYY-MM-DD   # When validation occurred
  validation_method: string     # How validation was done
  evidence:                     # Supporting evidence
    - type: string              # "field_data", "experiment", "natural_experiment"
      description: string
      source: string
  lessons_learned: string       # What we learned from validation
\end{verbatim}
\end{tcolorbox}


%%%%%%%%%%%%%%%%%%%%%%%%%%%%%%%%%%%%%%%%%%%%%%%%%%%%%%%%%%%%%%%%%%%%%%%%%%%%%%%
\section{Similarity Search}
\label{sec:fff-similarity}
%%%%%%%%%%%%%%%%%%%%%%%%%%%%%%%%%%%%%%%%%%%%%%%%%%%%%%%%%%%%%%%%%%%%%%%%%%%%%%%

The Registry must support efficient retrieval of similar models for the 3+1
Choice Architecture (Axiom MD-2 in Appendix UNMAPPED_DSN).

\subsection{Similarity Dimensions}

Models are compared across multiple dimensions:

\begin{table}[htbp]
\centering
\caption{Similarity Dimensions and Weights}
\label{tab:fff-similarity}
\renewcommand{\arraystretch}{1.3}
\begin{tabular}{@{}llcl@{}}
\toprule
\textbf{Dimension} & \textbf{Field(s)} & \textbf{Weight} & \textbf{Metric} \\
\midrule
Gestaltungsanlass & metadata.gestaltungsanlass\_type & 0.25 & Exact match \\
Domain & metadata.domain, metadata.tags & 0.20 & Jaccard similarity \\
Scope & input.scope.scope\_type & 0.20 & Exact match \\
Context & input.context\_psi.* & 0.20 & Weighted field match \\
Variables & specification.variables\_C & 0.10 & Set overlap \\
Validation & validation.status & 0.05 & Preference for validated \\
\bottomrule
\end{tabular}
\end{table}

\subsection{Similarity Algorithm}

\begin{tcolorbox}[colback=gray!5!white, colframe=gray!75!black,
    title=Axiom UNMAPPED_REG-1: Similarity Computation]
\textbf{Statement:} Given a query model $M_q$ and registry $\mathcal{R}$,
similarity is computed as:
\begin{equation}
\text{Sim}(M_q, M_i) = \sum_{d \in \mathcal{D}} w_d \cdot \text{sim}_d(M_q, M_i)
\end{equation}
where $\mathcal{D}$ is the set of dimensions, $w_d$ are weights (Table \ref{tab:fff-similarity}),
and $\text{sim}_d$ is the dimension-specific similarity metric.

\textbf{Top-K Selection:} The three most similar models are returned:
\begin{equation}
\{O_1, O_2, O_3\} = \text{TopK}(\{\text{Sim}(M_q, M_i) : M_i \in \mathcal{R}\}, k=3)
\end{equation}

\textbf{Diversity Constraint:} The three selected models must differ in at least
one material dimension to ensure meaningful choice (Axiom MD-1).
\end{tcolorbox}

\subsection{Cold Start Problem}

When the Registry is new or sparse:

\begin{tcolorbox}[colback=yellow!5!white, colframe=yellow!75!black,
    title=Axiom UNMAPPED_REG-2: Cold Start Handling]
\textbf{Statement:} If $|\mathcal{R}| < 3$ or fewer than 3 models exceed
similarity threshold $\tau = 0.3$, the workflow generates theory-derived
options instead:

\begin{equation}
\text{If } |\{M_i \in \mathcal{R} : \text{Sim}(M_q, M_i) > \tau\}| < 3:
\quad \text{Generate from EBF theory}
\end{equation}

\textbf{Marking:} Theory-derived options are clearly labeled as such, distinct
from registry-validated options.
\end{tcolorbox}


%%%%%%%%%%%%%%%%%%%%%%%%%%%%%%%%%%%%%%%%%%%%%%%%%%%%%%%%%%%%%%%%%%%%%%%%%%%%%%%
\section{Organizational Learning}
\label{sec:fff-learning}
%%%%%%%%%%%%%%%%%%%%%%%%%%%%%%%%%%%%%%%%%%%%%%%%%%%%%%%%%%%%%%%%%%%%%%%%%%%%%%%

The Registry enables organizational learning through systematic validation
feedback.

\subsection{Learning Loop}

\begin{figure}[htbp]
\centering
\begin{tikzpicture}[node distance=2cm, auto,
    block/.style={draw, rectangle, rounded corners, minimum width=3cm, minimum height=1cm, align=center}]

    % Nodes
    \node[block, fill=blue!10] (design) {Design Model\\(EEE Workflow)};
    \node[block, fill=green!10, right=of design] (register) {Register\\(FFF Schema)};
    \node[block, fill=yellow!10, right=of register] (apply) {Apply\\(Project)};
    \node[block, fill=orange!10, below=of apply] (observe) {Observe\\Outcomes};
    \node[block, fill=red!10, below=of register] (validate) {Validate\\Predictions};
    \node[block, fill=purple!10, below=of design] (learn) {Update\\Registry};

    % Arrows
    \draw[->] (design) -- (register);
    \draw[->] (register) -- (apply);
    \draw[->] (apply) -- (observe);
    \draw[->] (observe) -- (validate);
    \draw[->] (validate) -- (learn);
    \draw[->] (learn) -- (design);

    % Feedback to new designs
    \draw[->, dashed] (learn.west) -- ++(-1,0) |- (design.west);

\end{tikzpicture}
\caption{Organizational Learning Loop}
\label{fig:fff-learning}
\end{figure}

\subsection{Validation Protocol}

\begin{tcolorbox}[colback=gray!5!white, colframe=gray!75!black,
    title=Axiom UNMAPPED_REG-3: Mandatory Validation Review]
\textbf{Statement:} Every model with status ``pending'' must be reviewed for
validation within 12 months of application, or explicitly marked as
``validation\_not\_possible'' with documented reason.

\textbf{Validation Questions:}
\begin{enumerate}[nosep]
\item Were the predictions testable in practice?
\item What outcomes were observed?
\item Did predictions match outcomes within stated uncertainty?
\item What lessons does this provide for future models?
\end{enumerate}
\end{tcolorbox}

\subsection{Registry Metrics}

\begin{table}[htbp]
\centering
\caption{Registry Health Metrics}
\label{tab:fff-metrics}
\renewcommand{\arraystretch}{1.3}
\begin{tabular}{@{}lll@{}}
\toprule
\textbf{Metric} & \textbf{Definition} & \textbf{Target} \\
\midrule
Coverage & Domains with $\geq 5$ models & All active domains \\
Validation Rate & \% models validated within 12mo & $\geq 60\%$ \\
Accuracy & Mean accuracy of validated models & $\geq 0.7$ \\
Reuse Rate & \% new models using registry options & $\geq 50\%$ \\
Freshness & Mean age of top-retrieved models & $\leq 3$ years \\
\bottomrule
\end{tabular}
\end{table}


%%%%%%%%%%%%%%%%%%%%%%%%%%%%%%%%%%%%%%%%%%%%%%%%%%%%%%%%%%%%%%%%%%%%%%%%%%%%%%%
\section{Audit Trail}
\label{sec:fff-audit}
%%%%%%%%%%%%%%%%%%%%%%%%%%%%%%%%%%%%%%%%%%%%%%%%%%%%%%%%%%%%%%%%%%%%%%%%%%%%%%%

The Registry maintains complete audit trails for accountability.

\subsection{Traceability Requirements}

\begin{tcolorbox}[colback=gray!5!white, colframe=gray!75!black,
    title=Axiom UNMAPPED_REG-4: Complete Traceability]
\textbf{Statement:} Every client-facing statement or recommendation must be
traceable to:
\begin{enumerate}[nosep]
\item The specific model (Model\_ID)
\item The specific prediction(s) supporting it
\item The parameter values and sources used
\item The validation status at time of statement
\end{enumerate}

\textbf{Audit Query:} Given a statement, it must be possible to retrieve:
\begin{equation}
\text{Statement} \rightarrow \text{Model} \rightarrow \text{Predictions} \rightarrow \text{Parameters} \rightarrow \text{Sources}
\end{equation}
\end{tcolorbox}

\subsection{Version Control}

Models may be updated over time. Version control ensures:
\begin{itemize}
\item Historical versions remain accessible
\item Changes are documented with rationale
\item Statements reference specific versions
\end{itemize}


%%%%%%%%%%%%%%%%%%%%%%%%%%%%%%%%%%%%%%%%%%%%%%%%%%%%%%%%%%%%%%%%%%%%%%%%%%%%%%%

%%%%%%%%%%%%%%%%%%%%%%%%%%%%%%%%%%%%%%%%%%%%%%%%%%%%%%%%%%%%%%%%%%%%%%%%%%%%%%%
\section{Key Results}
\label{sec:results}
%%%%%%%%%%%%%%%%%%%%%%%%%%%%%%%%%%%%%%%%%%%%%%%%%%%%%%%%%%%%%%%%%%%%%%%%%%%%%%%

The analysis yields the following key results:

\begin{enumerate}
    \item \textbf{Result 1:} [TODO: Describe main finding]
    \item \textbf{Result 2:} [TODO: Describe secondary finding]
    \item \textbf{Result 3:} [TODO: Describe implication]
\end{enumerate}

\section{Integration with Other Appendices}
\label{sec:fff-integration}
%%%%%%%%%%%%%%%%%%%%%%%%%%%%%%%%%%%%%%%%%%%%%%%%%%%%%%%%%%%%%%%%%%%%%%%%%%%%%%%

\begin{tcolorbox}[colback=yellow!5!white, colframe=yellow!75!black,
    title=Integration: METHOD-REGISTRY $\leftrightarrow$ METHOD-DESIGN (EEE)]
\textbf{What METHOD-REGISTRY provides to EEE:}
\begin{itemize}[nosep]
\item Historical models for 3+1 option generation
\item Similarity search functionality
\item Validation status for option ranking
\end{itemize}

\textbf{What METHOD-REGISTRY requires from EEE:}
\begin{itemize}[nosep]
\item Completed models with full metadata
\item Adherence to schema structure
\item Quality standards for registration
\end{itemize}
\end{tcolorbox}

\begin{tcolorbox}[colback=yellow!5!white, colframe=yellow!75!black,
    title=Integration: METHOD-REGISTRY $\leftrightarrow$ DOMAIN Appendices]
\textbf{What METHOD-REGISTRY provides:}
\begin{itemize}[nosep]
\item Domain-specific model retrieval
\item Cross-domain learning opportunities
\end{itemize}

\textbf{What METHOD-REGISTRY requires:}
\begin{itemize}[nosep]
\item Domain tagging for all models
\item Domain-specific validation criteria
\end{itemize}
\end{tcolorbox}


% =============================================================================
%                    PART 3: APPLICATION
% =============================================================================

%%%%%%%%%%%%%%%%%%%%%%%%%%%%%%%%%%%%%%%%%%%%%%%%%%%%%%%%%%%%%%%%%%%%%%%%%%%%%%%
\section{Worked Example}
\label{sec:fff-example}
%%%%%%%%%%%%%%%%%%%%%%%%%%%%%%%%%%%%%%%%%%%%%%%%%%%%%%%%%%%%%%%%%%%%%%%%%%%%%%%

\subsection{Example: Complete Model Entry}

\begin{tcolorbox}[colback=white, colframe=black,
    title=Example Registry Entry: UNMAPPED_SHL-2026-001]
\begin{verbatim}
identification:
  model_id: UNMAPPED_SHL-2026-001
  designer: "Thomas Müller"
  created: 2026-01-15
  client: "PensionCo_CH"
  project: "Voluntary Savings Enhancement"
  version: 1.0

metadata:
  gestaltungsanlass: "Increase voluntary pension enrollment from 40% to 60%"
  gestaltungsanlass_type: opportunity
  entry_point: practice_driven
  domain: "pension"
  tags: ["retirement", "enrollment", "Switzerland", "B2C"]

input:
  scope:
    description: "Binary enrollment decision for voluntary pension savings"
    scope_type: decision
  context_psi:
    geographic: "Switzerland"
    demographic: "Employees 25-65, all income levels"
    temporal: "Decision at onboarding and annual review"
    channel: "Digital (app, email) + HR personal"
    regulatory: "Swiss BVG framework"

specification:
  variables_C: [F, E, S]
  complementarity_gamma:
    interactions: ["F*E", "temporal"]
    values: {"gamma_FE": 0.3, "gamma_temporal": 0.15}
  parameters_theta:
    values: {"beta_F": 0.45, "beta_E": 0.30, "beta_S": 0.25}
    sources: {"beta_F": "Thaler & Benartzi 2004", "beta_E": "Expert elicitation"}
    confidence: {"beta_F": 0.8, "beta_E": 0.6, "beta_S": 0.6}
  functional_form:
    type: "logistic"
    equation: "P(enroll) = 1/(1+exp(-(b0 + bF*F + bE*E + bS*S + gFE*F*E)))"

output:
  predictions:
    - id: P1
      statement: "Framing as 'future security' increases enrollment 8-12pp"
      type: comparative
      quantification: "8-12 percentage points"
    - id: P2
      statement: "Social proof increases enrollment 5-8pp"
      type: comparative
      quantification: "5-8 percentage points"
  statements:
    - date: 2026-02-01
      audience: "PensionCo Board"
      content: "Redesigned communication expected to increase participation by 15-20pp"
      based_on: [P1, P2]
  recommendations:
    - id: R1
      action: "Implement 'future security' framing in all communications"
      expected_impact: "+8-12pp enrollment"
      based_on: [P1]

validation:
  status: pending
  validation_date: null
  evidence: []
  lessons_learned: null
\end{verbatim}
\end{tcolorbox}


%%%%%%%%%%%%%%%%%%%%%%%%%%%%%%%%%%%%%%%%%%%%%%%%%%%%%%%%%%%%%%%%%%%%%%%%%%%%%%%
\section{Implications for Practice}
\label{sec:fff-practice}
%%%%%%%%%%%%%%%%%%%%%%%%%%%%%%%%%%%%%%%%%%%%%%%%%%%%%%%%%%%%%%%%%%%%%%%%%%%%%%%

\begin{tcolorbox}[colback=green!5!white, colframe=green!75!black,
    title=Practical Implications]
\begin{enumerate}
\item \textbf{Register every model:} Even ``small'' or ``quick'' models should
    be registered. Organizational memory depends on completeness.
\item \textbf{Complete all fields:} Partial entries reduce searchability and
    audit capability. Use ``unknown'' explicitly rather than leaving blank.
\item \textbf{Validate systematically:} Schedule validation reviews at project
    close or 12-month mark. Unvalidated models have limited learning value.
\item \textbf{Update on outcomes:} When outcomes become known, update the
    Registry promptly. This is how organizational learning happens.
\item \textbf{Use similarity search:} Before designing a new model, always
    search the Registry. Standing on shoulders beats reinventing wheels.
\end{enumerate}
\end{tcolorbox}


%%%%%%%%%%%%%%%%%%%%%%%%%%%%%%%%%%%%%%%%%%%%%%%%%%%%%%%%%%%%%%%%%%%%%%%%%%%%%%%
\section{Seed Models: DACH Reference Repository}
\label{sec:fff-seed}
%%%%%%%%%%%%%%%%%%%%%%%%%%%%%%%%%%%%%%%%%%%%%%%%%%%%%%%%%%%%%%%%%%%%%%%%%%%%%%%

The Registry is pre-populated with 8 validated reference models covering the
DACH region (Switzerland, Austria, Germany). These seed models serve as:
\begin{itemize}[nosep]
\item Starting points for new model design (3+1 options)
\item Templates demonstrating proper schema usage
\item Benchmarks for model quality
\end{itemize}

\begin{table}[htbp]
\centering
\caption{DACH Seed Models Overview}
\label{tab:fff-seed-overview}
\renewcommand{\arraystretch}{1.2}
\small
\begin{tabular}{@{}llllll@{}}
\toprule
\textbf{ID} & \textbf{Level} & \textbf{Geo} & \textbf{Domain} & \textbf{Scope} & \textbf{Use Case} \\
\midrule
SEED-CH-01 & L=0 & 🇨🇭 & Vorsorge & Decision & PK Enrollment \\
SEED-CH-02 & L=0 & 🇨🇭 & Vorsorge & Continuous & Beitragshöhe \\
SEED-CH-03 & L=3 & 🇨🇭 & HR/Org & Process & Change Adoption \\
SEED-CH-04 & L=1 & 🇨🇭 & Energy & Continuous & Verbrauchsreduktion \\
\addlinespace
SEED-UNMAPPED_P06-01 & L=0 & 🇦🇹 & Gesundheit & Decision & Vorsorgeuntersuchung \\
SEED-UNMAPPED_P06-02 & L=0 & 🇦🇹 & Public & Decision & Steuer Compliance \\
\addlinespace
SEED-DE-01 & L=1 & 🇩🇪 & Finance & Decision & Riester Familie \\
SEED-DE-02 & L=3 & 🇩🇪 & HR/Org & Process & Digitalisierung \\
\bottomrule
\end{tabular}
\end{table}

% ---------------------------------------------------------------------------
\subsection{Switzerland: SEED-CH-01 -- Pensionskasse Enrollment}
% ---------------------------------------------------------------------------

\begin{tcolorbox}[colback=white, colframe=red!50!black,
    title=SEED-CH-01: Pensionskasse Voluntary Enrollment (Switzerland)]
\begin{verbatim}
identification:
  model_id: SEED-CH-01
  designer: "FehrAdvice Reference Team"
  created: 2026-01-01
  client: "Reference Model"
  project: "DACH Seed Repository"
  version: 1.0

metadata:
  gestaltungsanlass: "Increase voluntary pension (PK) enrollment rate"
  gestaltungsanlass_type: opportunity
  entry_point: practice_driven
  domain: "vorsorge"
  tags: ["pension", "enrollment", "Switzerland", "B2C", "BVG"]

input:
  scope:
    description: "Binary enrollment decision for voluntary PK contributions"
    scope_type: decision
  context_psi:
    geographic: "Switzerland"
    demographic: "Employees 25-65, income CHF 60k+"
    temporal: "Decision at onboarding and annual review"
    channel: "Digital (HR portal) + HR personal consultation"
    regulatory: "BVG/BVV2 framework, tax-deductible"
    cultural: "High autonomy, strong savings culture"
  constraints:
    data_available: ["enrollment_rates", "demographics", "channel_usage"]

specification:
  variables_C: [F, E, S]
  complementarity_gamma:
    interactions: ["F*E", "temporal"]
    values:
      gamma_FE: 0.30
      gamma_temporal: 0.20
  parameters_theta:
    values:
      beta_F: 0.35
      beta_E: 0.25
      beta_S: 0.20
      beta_0: -1.50
    sources:
      beta_F: "Thaler & Benartzi 2004, Swiss adjustment"
      beta_E: "Expert elicitation FehrAdvice"
      beta_S: "Duflo & Saez 2003"
    confidence:
      beta_F: 0.75
      beta_E: 0.60
      beta_S: 0.70
  functional_form:
    type: "logistic"
    equation: "P(enroll) = 1/(1+exp(-(b0 + bF*F + bE*E + bS*S + gFE*F*E)))"
    justification: "Binary outcome, standard choice model"

output:
  predictions:
    - id: P1
      statement: "Framing contributions as 'future security' increases
                  enrollment by 8-12 percentage points"
      type: comparative
      quantification: "8-12pp"
      testable: true
    - id: P2
      statement: "Social proof ('X% of colleagues enrolled') increases
                  enrollment by 5-8 percentage points"
      type: comparative
      quantification: "5-8pp"
      testable: true
    - id: P3
      statement: "Effect of P1 and P2 is superadditive due to gamma_FE > 0"
      type: comparative
      quantification: "Combined effect 15-22pp vs. sum 13-20pp"
      testable: true

validation:
  status: validated
  accuracy:
    overall: 0.78
    by_prediction: {P1: 0.82, P2: 0.75, P3: 0.68}
  validation_date: 2025-06-15
  validation_method: "Field experiment with 3 PK providers, n=4,200"
  evidence:
    - type: "field_experiment"
      description: "RCT with randomized framing conditions"
      source: "FehrAdvice internal study FA-2025-PK-01"
  lessons_learned: "gamma_FE confirmed; social proof effect varies by age"
\end{verbatim}
\end{tcolorbox}

% ---------------------------------------------------------------------------
\subsection{Switzerland: SEED-CH-02 -- Contribution Level}
% ---------------------------------------------------------------------------

\begin{tcolorbox}[colback=white, colframe=red!50!black,
    title=SEED-CH-02: Voluntary Contribution Amount (Switzerland)]
\begin{verbatim}
identification:
  model_id: SEED-CH-02
  designer: "FehrAdvice Reference Team"
  created: 2026-01-01
  client: "Reference Model"
  project: "DACH Seed Repository"
  version: 1.0

metadata:
  gestaltungsanlass: "Increase average voluntary contribution amount"
  gestaltungsanlass_type: opportunity
  entry_point: practice_driven
  domain: "vorsorge"
  tags: ["pension", "contribution", "Switzerland", "B2C", "continuous"]

input:
  scope:
    description: "Continuous choice of monthly contribution amount (CHF)"
    scope_type: continuous
  context_psi:
    geographic: "Switzerland"
    demographic: "Employees already enrolled, income CHF 80k+"
    temporal: "Annual adjustment window"
    channel: "Digital self-service"
    regulatory: "BVG max contribution limits"

specification:
  variables_C: [F, E, S, D]
  complementarity_gamma:
    interactions: ["F*E", "F*D", "temporal"]
    values:
      gamma_FE: 0.25
      gamma_FD: 0.15
      gamma_temporal: 0.30
  parameters_theta:
    values:
      beta_F: 0.40
      beta_E: 0.20
      beta_S: 0.15
      beta_D: 0.10
    sources:
      beta_F: "Madrian & Shea 2001"
      beta_D: "Hershfield et al. 2011 (future self)"
  functional_form:
    type: "linear_with_anchoring"
    equation: "Y = anchor + beta_F*F + beta_E*E + beta_S*S + beta_D*D + eps"
    justification: "Continuous outcome, strong anchoring effects"

output:
  predictions:
    - id: P1
      statement: "Showing 'recommended' contribution increases average by 15-20%"
      type: comparative
      quantification: "15-20% increase"
    - id: P2
      statement: "Future-self visualization increases contribution by 10-15%"
      type: comparative
      quantification: "10-15% increase"

validation:
  status: validated
  accuracy:
    overall: 0.72
  validation_date: 2025-09-01
  validation_method: "A/B test with 2 PK providers"
\end{verbatim}
\end{tcolorbox}

% ---------------------------------------------------------------------------
\subsection{Switzerland: SEED-CH-03 -- Organizational Change Adoption}
% ---------------------------------------------------------------------------

\begin{tcolorbox}[colback=white, colframe=red!50!black,
    title=SEED-CH-03: Change Adoption in Swiss Organizations (L=3)]
\begin{verbatim}
identification:
  model_id: SEED-CH-03
  designer: "FehrAdvice Reference Team"
  created: 2026-01-01
  client: "Reference Model"
  project: "DACH Seed Repository"
  version: 1.0

metadata:
  gestaltungsanlass: "Accelerate adoption of new processes/tools in organization"
  gestaltungsanlass_type: situation
  entry_point: practice_driven
  domain: "hr_org"
  tags: ["change", "adoption", "Switzerland", "B2B", "process", "L3"]

input:
  scope:
    description: "Organizational journey through change adoption stages"
    scope_type: process
  context_psi:
    geographic: "Switzerland"
    demographic: "Mid-size companies, 100-500 employees"
    temporal: "6-18 month change programs"
    channel: "Internal communication + training"
    regulatory: "Swiss labor law consultation requirements"
    cultural: "Consensus-oriented, high employee voice"

specification:
  variables_C: [F, E, S, D]
  complementarity_gamma:
    interactions: ["S*E", "D*F", "cross_L0_L3"]
    values:
      gamma_SE: 0.35
      gamma_DF: 0.20
      gamma_cross: 0.25
  parameters_theta:
    values:
      beta_S: 0.30
      beta_E: 0.25
      beta_D: 0.20
      beta_F: 0.15
      tau: 90  # days time constant
    sources:
      gamma_SE: "Organizational behavior literature"
      tau: "FehrAdvice change projects 2020-2025"
  functional_form:
    type: "bcj_dynamic"
    equation: "dS/dt = (1/tau) * [S*(A, WAX, Psi) - S(t)] + beta_social * S_peers"
    justification: "Process scope requires BCJ dynamics"

output:
  predictions:
    - id: P1
      statement: "Leadership visible commitment reduces tau by 30-40%"
      type: comparative
      quantification: "tau reduction 30-40%"
    - id: P2
      statement: "Peer adoption above 30% triggers acceleration (social tipping)"
      type: boundary
      quantification: "Tipping point at ~30% adoption"
    - id: P3
      statement: "Cross-level alignment (L0-L3) increases final adoption by 20pp"
      type: comparative
      quantification: "+20pp final adoption"

validation:
  status: partially_validated
  accuracy:
    overall: 0.65
    by_prediction: {P1: 0.70, P2: 0.72, P3: 0.55}
  validation_date: 2025-11-01
  validation_method: "Longitudinal study across 8 change projects"
  lessons_learned: "P3 (cross-level) needs more data; cultural factors stronger than expected"
\end{verbatim}
\end{tcolorbox}

% ---------------------------------------------------------------------------
\subsection{Switzerland: SEED-CH-04 -- Energy Consumption Reduction}
% ---------------------------------------------------------------------------

\begin{tcolorbox}[colback=white, colframe=red!50!black,
    title=SEED-CH-04: Household Energy Reduction (Switzerland, L=1)]
\begin{verbatim}
identification:
  model_id: SEED-CH-04
  designer: "FehrAdvice Reference Team"
  created: 2026-01-01
  client: "Reference Model"
  project: "DACH Seed Repository"
  version: 1.0

metadata:
  gestaltungsanlass: "Reduce household electricity consumption"
  gestaltungsanlass_type: opportunity
  entry_point: hybrid
  domain: "energy"
  tags: ["energy", "sustainability", "Switzerland", "household", "L1"]

input:
  scope:
    description: "Continuous reduction in monthly kWh consumption"
    scope_type: continuous
  context_psi:
    geographic: "Switzerland"
    demographic: "Households, owners and renters"
    temporal: "Monthly measurement, seasonal variation"
    channel: "Smart meter app + utility communication"
    regulatory: "Energiestrategie 2050"
    cultural: "High environmental awareness"

specification:
  variables_C: [F, E, S, X]
  complementarity_gamma:
    interactions: ["F*X", "S*X", "temporal"]
    values:
      gamma_FX: 0.20
      gamma_SX: 0.30
      gamma_temporal: 0.25
  parameters_theta:
    values:
      beta_F: 0.35
      beta_E: 0.15
      beta_S: 0.20
      beta_X: 0.15
    sources:
      beta_S: "Allcott 2011 (OPOWER)"
      beta_X: "Swiss Energy Survey 2024"
  functional_form:
    type: "linear_percentage"
    equation: "Y_reduction = beta_F*F + beta_S*S + beta_X*X + gamma_SX*S*X + eps"

output:
  predictions:
    - id: P1
      statement: "Social comparison feedback reduces consumption 5-8%"
      type: comparative
      quantification: "5-8% reduction"
    - id: P2
      statement: "Environmental framing + social proof is superadditive (gamma_SX)"
      type: comparative
      quantification: "Combined 10-14% vs. sum 8-12%"
    - id: P3
      statement: "Effect decays without reinforcement (tau ~ 60 days)"
      type: boundary
      quantification: "Half-life ~60 days"

validation:
  status: validated
  accuracy:
    overall: 0.75
  validation_date: 2025-08-01
  validation_method: "Utility partnership pilot, n=12,000 households"
\end{verbatim}
\end{tcolorbox}

% ---------------------------------------------------------------------------
\subsection{Austria: SEED-UNMAPPED_P06-01 -- Preventive Health Screening}
% ---------------------------------------------------------------------------

\begin{tcolorbox}[colback=white, colframe=blue!50!black,
    title=SEED-UNMAPPED_P06-01: Vorsorgeuntersuchung Participation (Austria)]
\begin{verbatim}
identification:
  model_id: SEED-UNMAPPED_P06-01
  designer: "FehrAdvice Reference Team"
  created: 2026-01-01
  client: "Reference Model"
  project: "DACH Seed Repository"
  version: 1.0

metadata:
  gestaltungsanlass: "Increase participation in preventive health screenings"
  gestaltungsanlass_type: opportunity
  entry_point: practice_driven
  domain: "gesundheit"
  tags: ["health", "prevention", "Austria", "B2C", "ASVG"]

input:
  scope:
    description: "Binary decision to attend annual Vorsorgeuntersuchung"
    scope_type: decision
  context_psi:
    geographic: "Austria"
    demographic: "Adults 18+, ASVG-insured"
    temporal: "Annual reminder cycle"
    channel: "GKK letter + digital reminder"
    regulatory: "ASVG preventive care coverage"
    cultural: "Moderate health consciousness, trust in state system"

specification:
  variables_C: [P, E, F, S]
  complementarity_gamma:
    interactions: ["P*E", "temporal"]
    values:
      gamma_PE: 0.40
      gamma_temporal: 0.35
  parameters_theta:
    values:
      beta_P: 0.35
      beta_E: 0.30
      beta_F: 0.10
      beta_S: 0.15
    sources:
      gamma_temporal: "Present bias literature (O'Donoghue & Rabin)"
  functional_form:
    type: "logistic"
    equation: "P(attend) = 1/(1+exp(-(b0 + bP*P + bE*E + gPE*P*E)))"

output:
  predictions:
    - id: P1
      statement: "Implementation intentions ('when/where') increase attendance 12-18pp"
      type: comparative
      quantification: "12-18pp"
    - id: P2
      statement: "Loss-framed health messages outperform gain-framed by 5-8pp"
      type: comparative
      quantification: "5-8pp advantage"
    - id: P3
      statement: "SMS reminder 48h before reduces no-shows by 25-35%"
      type: comparative
      quantification: "25-35% no-show reduction"

validation:
  status: validated
  accuracy:
    overall: 0.70
  validation_date: 2025-05-01
  validation_method: "GKK Oberösterreich pilot program"
\end{verbatim}
\end{tcolorbox}

% ---------------------------------------------------------------------------
\subsection{Austria: SEED-UNMAPPED_P06-02 -- Tax Compliance}
% ---------------------------------------------------------------------------

\begin{tcolorbox}[colback=white, colframe=blue!50!black,
    title=SEED-UNMAPPED_P06-02: Tax Declaration Timeliness (Austria)]
\begin{verbatim}
identification:
  model_id: SEED-UNMAPPED_P06-02
  designer: "FehrAdvice Reference Team"
  created: 2026-01-01
  client: "Reference Model"
  project: "DACH Seed Repository"
  version: 1.0

metadata:
  gestaltungsanlass: "Increase on-time tax declaration submission"
  gestaltungsanlass_type: problem
  entry_point: practice_driven
  domain: "public"
  tags: ["tax", "compliance", "Austria", "public sector", "nudging"]

input:
  scope:
    description: "Binary: on-time vs. late submission of Einkommensteuererklärung"
    scope_type: decision
  context_psi:
    geographic: "Austria"
    demographic: "Self-employed and voluntary filers"
    temporal: "Annual deadline April 30 / June 30 (FinanzOnline)"
    channel: "FinanzOnline + postal reminder"
    regulatory: "BAO deadlines, penalty structure"

specification:
  variables_C: [F, S, E]
  complementarity_gamma:
    interactions: ["S*F", "trust"]
    values:
      gamma_SF: 0.20
      gamma_trust: 0.25
  parameters_theta:
    values:
      beta_F: 0.30
      beta_S: 0.25
      beta_E: 0.15
    sources:
      beta_S: "Hallsworth et al. 2017 (UK tax letters)"
  functional_form:
    type: "logistic"

output:
  predictions:
    - id: P1
      statement: "Social norm message ('9 out of 10 pay on time') increases
                  compliance by 3-5pp"
      type: comparative
      quantification: "3-5pp"
    - id: P2
      statement: "Simplified deadline reminder increases compliance by 2-4pp"
      type: comparative
      quantification: "2-4pp"
    - id: P3
      statement: "Combination with late fee salience increases effect to 6-9pp"
      type: comparative
      quantification: "6-9pp combined"

validation:
  status: validated
  accuracy:
    overall: 0.73
  validation_date: 2025-07-01
  validation_method: "BMF Austria pilot, n=50,000 letters"
\end{verbatim}
\end{tcolorbox}

% ---------------------------------------------------------------------------
\subsection{Germany: SEED-DE-01 -- Riester Family Decision}
% ---------------------------------------------------------------------------

\begin{tcolorbox}[colback=white, colframe=yellow!50!black,
    title=SEED-DE-01: Riester Pension Family Decision (Germany, L=1)]
\begin{verbatim}
identification:
  model_id: SEED-DE-01
  designer: "FehrAdvice Reference Team"
  created: 2026-01-01
  client: "Reference Model"
  project: "DACH Seed Repository"
  version: 1.0

metadata:
  gestaltungsanlass: "Increase Riester pension uptake among eligible families"
  gestaltungsanlass_type: opportunity
  entry_point: practice_driven
  domain: "finance"
  tags: ["pension", "Riester", "Germany", "family", "L1", "Kinderzulage"]

input:
  scope:
    description: "Household decision to open/maintain Riester contract"
    scope_type: decision
  context_psi:
    geographic: "Germany"
    demographic: "Families with children, income EUR 30-70k"
    temporal: "Year-end decision, Zulageantrag deadline"
    channel: "Bank/insurance advisor + digital"
    regulatory: "Altersvermögensgesetz, Kinderzulage EUR 300/child"
    cultural: "State-reliance orientation, complex product perception"

specification:
  variables_C: [F, E, S]
  complementarity_gamma:
    interactions: ["F*S", "complexity_barrier"]
    values:
      gamma_FS: 0.25
      gamma_complexity: -0.30
  parameters_theta:
    values:
      beta_F: 0.45
      beta_E: 0.20
      beta_S: 0.25
    sources:
      beta_F: "Riester uptake studies (Börsch-Supan)"
      gamma_complexity: "DIA Studie 2024"
  functional_form:
    type: "logistic_with_friction"
    equation: "P(riester) = 1/(1+exp(-(b0 + bF*F - friction + bS*S)))"

output:
  predictions:
    - id: P1
      statement: "Simplified calculator showing Kinderzulage increases uptake 8-12pp"
      type: comparative
      quantification: "8-12pp"
    - id: P2
      statement: "Pre-filled Zulageantrag reduces dropout by 15-25pp"
      type: comparative
      quantification: "15-25pp retention"
    - id: P3
      statement: "Social proof ('Familien wie Ihre') adds 4-6pp"
      type: comparative
      quantification: "4-6pp"

validation:
  status: partially_validated
  accuracy:
    overall: 0.62
  validation_date: 2025-10-01
  validation_method: "Insurance partner pilot, n=8,000"
  lessons_learned: "Complexity barrier (-0.30) confirmed; simplification critical"
\end{verbatim}
\end{tcolorbox}

% ---------------------------------------------------------------------------
\subsection{Germany: SEED-DE-02 -- Digitalization Adoption}
% ---------------------------------------------------------------------------

\begin{tcolorbox}[colback=white, colframe=yellow!50!black,
    title=SEED-DE-02: Digitalization Adoption in German SME (L=3)]
\begin{verbatim}
identification:
  model_id: SEED-DE-02
  designer: "FehrAdvice Reference Team"
  created: 2026-01-01
  client: "Reference Model"
  project: "DACH Seed Repository"
  version: 1.0

metadata:
  gestaltungsanlass: "Accelerate digital tool adoption in German Mittelstand"
  gestaltungsanlass_type: situation
  entry_point: practice_driven
  domain: "hr_org"
  tags: ["digitalization", "SME", "Germany", "Mittelstand", "L3", "process"]

input:
  scope:
    description: "Organizational journey through digital transformation stages"
    scope_type: process
  context_psi:
    geographic: "Germany"
    demographic: "SME/Mittelstand, 50-250 employees, manufacturing focus"
    temporal: "12-24 month transformation programs"
    channel: "Management + Betriebsrat + training"
    regulatory: "Betriebsverfassungsgesetz, Mitbestimmung"
    cultural: "Hierarchical, risk-averse, process-oriented"

specification:
  variables_C: [F, E, S, D]
  complementarity_gamma:
    interactions: ["S*E", "D*F", "cross_L0_L3", "betriebsrat"]
    values:
      gamma_SE: 0.30
      gamma_DF: 0.25
      gamma_cross: 0.20
      gamma_BR: 0.35
  parameters_theta:
    values:
      beta_S: 0.25
      beta_E: 0.30
      beta_D: 0.15
      beta_F: 0.20
      tau: 120  # days, longer than CH due to process orientation
    sources:
      tau: "German transformation case studies"
      gamma_BR: "Mitbestimmung literature"
  functional_form:
    type: "bcj_dynamic"
    equation: "dS/dt = (1/tau) * [S*(A, WAX, Psi) - S(t)] + gamma_BR * BR_support"

output:
  predictions:
    - id: P1
      statement: "Betriebsrat early involvement reduces tau by 25-35%"
      type: comparative
      quantification: "tau reduction 25-35%"
    - id: P2
      statement: "Qualification guarantee ('no job loss') increases adoption 15-20pp"
      type: comparative
      quantification: "15-20pp"
    - id: P3
      statement: "Pilot group success stories accelerate diffusion (social proof)"
      type: comparative
      quantification: "Diffusion speed +40%"

validation:
  status: partially_validated
  accuracy:
    overall: 0.58
  validation_date: 2025-12-01
  validation_method: "4 Mittelstand digitalization projects"
  lessons_learned: "Betriebsrat effect (gamma_BR) stronger than expected;
                    hierarchy slows horizontal diffusion"
\end{verbatim}
\end{tcolorbox}

\subsection{Seed Model Usage Guidelines}

\begin{tcolorbox}[colback=green!5!white, colframe=green!75!black,
    title=How to Use Seed Models]
\begin{enumerate}
\item \textbf{As 3+1 Options:} When designing a new model, similarity search
    returns relevant seed models as starting options.
\item \textbf{As Templates:} Copy schema structure; replace parameters with
    project-specific values.
\item \textbf{As Benchmarks:} Compare new model predictions against seed model
    accuracy to assess quality.
\item \textbf{Update When Validated:} If your project validates/improves a seed
    model, propose updates to the reference repository.
\end{enumerate}

\textbf{Important:} Seed models are \textit{starting points}, not final answers.
Always adjust for project-specific context ($\Psi$) and validate predictions.
\end{tcolorbox}


% =============================================================================
%                    PART 4: BACK MATTER
% =============================================================================

%%%%%%%%%%%%%%%%%%%%%%%%%%%%%%%%%%%%%%%%%%%%%%%%%%%%%%%%%%%%%%%%%%%%%%%%%%%%%%%
\section{Summary}
\label{sec:fff-summary}
%%%%%%%%%%%%%%%%%%%%%%%%%%%%%%%%%%%%%%%%%%%%%%%%%%%%%%%%%%%%%%%%%%%%%%%%%%%%%%%

\begin{tcolorbox}[colback=green!10!white, colframe=green!75!black,
    title=\textbf{Appendix UNMAPPED_REG: Summary}]

\textbf{Core Question:} How do we store, retrieve, and learn from behavioral models?

\textbf{Main Contribution:}
\begin{itemize}[nosep]
    \item Defines the 6-component Model Schema
    \item Specifies similarity search for 3+1 option retrieval
    \item Establishes organizational learning loop
    \item Ensures audit trail for recommendations
\end{itemize}

\textbf{Key Axioms:}
\begin{itemize}[nosep]
    \item UNMAPPED_REG-1: Similarity Computation
    \item UNMAPPED_REG-2: Cold Start Handling
    \item UNMAPPED_REG-3: Mandatory Validation Review
    \item UNMAPPED_REG-4: Complete Traceability
\end{itemize}

\textbf{Integration:} This appendix connects to EEE (Design Workflow),
BBB (Parameters), and all DOMAIN/PREDICT appendices.

\end{tcolorbox}


%%%%%%%%%%%%%%%%%%%%%%%%%%%%%%%%%%%%%%%%%%%%%%%%%%%%%%%%%%%%%%%%%%%%%%%%%%%%%%%
\section{Glossary of Symbols}
\label{sec:fff-glossary}
%%%%%%%%%%%%%%%%%%%%%%%%%%%%%%%%%%%%%%%%%%%%%%%%%%%%%%%%%%%%%%%%%%%%%%%%%%%%%%%

\begin{tcolorbox}[colback=blue!5!white, colframe=blue!75!black]
\textbf{Central Reference:} For complete symbol definitions, see
\textbf{Appendix UNMAPPED_GLS} (\textit{Glossary and Notation}, \S\ref{app:glossary}).

This section lists only symbols introduced or specialized in this appendix.
\end{tcolorbox}

\vspace{0.5em}

\begin{table}[htbp]
\centering
\caption{Symbols Introduced in Appendix UNMAPPED_REG}
\label{tab:fff-symbols}
\renewcommand{\arraystretch}{1.3}
\begin{tabular}{@{}clp{6cm}c@{}}
\toprule
\textbf{Symbol} & \textbf{Name} & \textbf{Definition} & \textbf{Also in G?} \\
\midrule
$\mathcal{R}$ & Registry & Set of all stored models & NEW \\
$M_i$ & Model & A complete model entry & \checkmark \\
$\text{Sim}(\cdot,\cdot)$ & Similarity & Similarity function between models & NEW \\
$w_d$ & Weight & Dimension weight in similarity & NEW \\
$\tau$ & Threshold & Minimum similarity for retrieval & NEW \\
$\mathcal{D}$ & Dimensions & Set of comparison dimensions & NEW \\
\bottomrule
\end{tabular}
\end{table}


%%%%%%%%%%%%%%%%%%%%%%%%%%%%%%%%%%%%%%%%%%%%%%%%%%%%%%%%%%%%%%%%%%%%%%%%%%%%%%%

%%%%%%%%%%%%%%%%%%%%%%%%%%%%%%%%%%%%%%%%%%%%%%%%%%%%%%%%%%%%%%%%%%%%%%%%%%%%%%%
\section{Critical Foundations and Objections}
\label{sec:reg-foundations}
%%%%%%%%%%%%%%%%%%%%%%%%%%%%%%%%%%%%%%%%%%%%%%%%%%%%%%%%%%%%%%%%%%%%%%%%%%%%%%%

\subsection{Objection 1: Scope Limitations}

\textbf{Objection:} The framework presented may have limited applicability
outside the specific domain contexts discussed.

\textbf{Response:} We acknowledge domain-specificity. The framework provides
general principles that should be calibrated to specific contexts. Cross-domain
validation remains an open research question.

\subsection{Objection 2: Measurement Challenges}

\textbf{Objection:} Key constructs may be difficult to measure reliably in
practice.

\textbf{Response:} We recommend using validated scales and instruments where
available, and developing domain-specific proxies where necessary. Sensitivity
analysis should assess robustness to measurement uncertainty.


\section{Open Issues and Research Agenda}
\label{sec:fff-open}
%%%%%%%%%%%%%%%%%%%%%%%%%%%%%%%%%%%%%%%%%%%%%%%%%%%%%%%%%%%%%%%%%%%%%%%%%%%%%%%

\begin{table}[htbp]
\centering
\caption{Research Roadmap for METHOD-REGISTRY}
\label{tab:fff-roadmap}
\renewcommand{\arraystretch}{1.3}
\begin{tabular}{@{}clcl@{}}
\toprule
\textbf{\#} & \textbf{Issue} & \textbf{Priority} & \textbf{Status} \\
\midrule
1 & Optimal similarity weights & HIGH & Empirical calibration needed \\
2 & Automated validation detection & MEDIUM & Exploratory \\
3 & Cross-organizational registries & LOW & Future consideration \\
4 & Privacy-preserving model sharing & MEDIUM & Design phase \\
5 & ML-enhanced similarity search & LOW & Research \\
\bottomrule
\end{tabular}
\end{table}


%%%%%%%%%%%%%%%%%%%%%%%%%%%%%%%%%%%%%%%%%%%%%%%%%%%%%%%%%%%%%%%%%%%%%%%%%%%%%%%
\section{References}
\label{sec:fff-references}
%%%%%%%%%%%%%%%%%%%%%%%%%%%%%%%%%%%%%%%%%%%%%%%%%%%%%%%%%%%%%%%%%%%%%%%%%%%%%%%

\begin{tcolorbox}[colback=blue!5!white, colframe=blue!75!black]
\textbf{Master Bibliography:} For complete reference details, see
\texttt{bcm2\_master\_references.bib} in the repository.

This section lists appendix-specific references and data sources.
\end{tcolorbox}

\vspace{0.5em}

\subsection{Key References for This Appendix}

\begin{itemize}[nosep]
    \item \textbf{Knowledge Management:} Nonaka \& Takeuchi (1995), Davenport \& Prusak (1998)
    \item \textbf{Organizational Learning:} Argyris \& Schön (1978), Senge (1990)
    \item \textbf{Model Registries:} MLflow documentation, ModelDB (Vartak et al., 2016)
\end{itemize}

\nocite{bcm_master}


%%%%%%%%%%%%%%%%%%%%%%%%%%%%%%%%%%%%%%%%%%%%%%%%%%%%%%%%%%%%%%%%%%%%%%%%%%%%%%%
% CROSS-REFERENCE FOOTER (Required)
%%%%%%%%%%%%%%%%%%%%%%%%%%%%%%%%%%%%%%%%%%%%%%%%%%%%%%%%%%%%%%%%%%%%%%%%%%%%%%%

\vspace{1em}
\hrule
\vspace{0.5em}

\noindent\textit{Cross-references:}
Appendix UNMAPPED_GLS (Central Glossary),
Appendix UNMAPPED_DSN (METHOD-DESIGN),
Appendix UNMAPPED_EST (CORE-WHERE),
All DOMAIN- appendices.

\noindent\textit{Master Bibliography:} \texttt{bcm2\_master\_references.bib}


% =============================================================================
% END OF APPENDIX FFF
% =============================================================================
